\documentclass[11pt]{article}
\usepackage[margin=1in]{geometry}
\usepackage{amsmath,amssymb,amsthm}
\usepackage{booktabs}
\usepackage{array}
\usepackage[hidelinks]{hyperref}
\usepackage{xcolor}
\usepackage{listings}
\lstset{basicstyle=\ttfamily\small,columns=fullflexible,breaklines=true,frame=single,framerule=0.3pt}

\newtheorem{theorem}{Theorem}
\newtheorem{proposition}[theorem]{Proposition}
\newtheorem{lemma}[theorem]{Lemma}
\newtheorem{corollary}[theorem]{Corollary}
\theoremstyle{remark}
\newtheorem{remark}[theorem]{Remark}
\theoremstyle{definition}
\newtheorem{definition}[theorem]{Definition}

\newcommand{\Z}{\mathbb{Z}}
\newcommand{\Q}{\mathbb{Q}}
\newcommand{\leg}[2]{\left(\frac{#1}{#2}\right)}

\title{A large integral point on the cubic surface\\[2pt]
$x^2y+xy^2+xz^2+2x^2-z^2+3=0$}
\author{Computation and exposition: Claude (Anthropic)\\ prepared for Jamal Agbanwa}
\date{25 October 2026}

\begin{document}
\maketitle

\begin{abstract}
We report the integer solution
\[
(x,y,z)=(-2450412151997178,\ 2357619934960084,\ 467727250324161)
\]
of $x^2y+xy^2+xz^2+2x^2-z^2+3=0$, together with the three companions obtained from the
symmetries $z\mapsto -z$ and $y\mapsto -x-y$.  It was found by a complete, sieve-based search over
$x$ that ran for about two days on a 2-core cloud machine and an 8-core Apple M1 Pro.  We explain
the elementary reductions and the local--global arithmetic (including a Brauer--Manin type
reciprocity condition) that confine $x$ to $x\equiv 6 \pmod{72}$ and to a set of density
about $1/1800$ in the integers, why one should expect the
first solution to be large, why ``small-parameter'' methods (elliptic fibrations, caps) could
not have found this particular point, and how the search was organised, validated and
executed.  The search was completed up to $|x|<2.452\cdot10^{15}$.  Up to the symmetries above,
the point is the only solution in that range, and so the one with the smallest $|x|$
(Theorem~\ref{thm:unique}).
\end{abstract}

\tableofcontents

\section{The result}

\begin{theorem}\label{thm:main}
The quadruples
\[
\begin{array}{rcl}
x&=&-2450412151997178,\\
y&\in&\{\,2357619934960084,\ \ 92792217037094\,\},\\
z&=&\pm 467727250324161
\end{array}
\]
(four triples in total) satisfy
\begin{equation}\label{eq:main}
x^2y+xy^2+xz^2+2x^2-z^2+3=0 .
\end{equation}
\end{theorem}

\begin{proof}
Direct computation.  For
\[(x,y,z)=(-2450412151997178,\,2357619934960084,\,467727250324161)\]
the six terms are:
\[
\begin{array}{lr}
x^2y   & 14156375379132502896580275368792446028069589456\\
xy^2   & -13620301700437919819452512803799197104495927968\\
xz^2   & -536073678694594867398411180074838272233234938\\
2x^2   & 12009039429310881955628239927368\\
-z^2   & -218768780695800366279580353921\\
+3     & 3\\\midrule
\text{sum} & 0
\end{array}
\]
The left side of \eqref{eq:main} is even in $z$, and it is invariant under $y\mapsto -x-y$
because $x^2y+xy^2=xy(x+y)$ is.  This gives $-x-y=92792217037094$.
\end{proof}

The numbers have 15--16 decimal digits.  The rest of this note explains why a solution of this
size was to be expected, and how it was found.

\section{First reductions}

\subsection{Parity}
If $x$ is odd then $xy(x+y)$ is even, $z^2(x-1)$ is even and $2x^2$ is even, so the left side of
\eqref{eq:main} is odd.  Hence \emph{$x$ is even}.  Write
\[
x=2a,\qquad y=b-a\qquad (a,b\in\Z).
\]
Since $x^2y+xy^2=xy(x+y)=2a(b-a)(b+a)=2a(b^2-a^2)$, equation \eqref{eq:main} becomes
\begin{equation}\label{eq:red}
\boxed{\;2a\,b^2+(2a-1)\,z^2 \;=\; N(a):=2a^3-8a^2-3\;}
\end{equation}
For the solution of Theorem~\ref{thm:main},
\[a=-1225206075998589,\quad b=1132413858961495,\quad z=467727250324161.\]

\subsection{The conic bundle and the circle at infinity}
For fixed $a$, \eqref{eq:red} is a conic in $(b,z)$.  For $a>0$ both coefficients are positive.
For $a<0$ one multiplies by $-1$.  Either way it is an \emph{ellipse}, so each fibre carries
finitely many integral points.  Dividing by $2a$ gives
\[
b^2+z^2\Bigl(1-\frac1{2a}\Bigr)=a^2-4a-\frac{3}{2a},
\]
so asymptotically the point $(b,z)$ lies on the circle $b^2+z^2\approx a^2$.  It is natural to write
\[
b\approx |a|\cos\varphi,\qquad z\approx |a|\sin\varphi .
\]
The found solution has $\varphi=\arctan(z/b)\approx 0.3917$ rad ($22.4^\circ$), with
$b/|a|\approx 0.924$ and $z/|a|\approx 0.382$.  This is a perfectly ``generic'' angle, a fact that
matters in Section~\ref{sec:caps}.

A second useful quantity is
\begin{equation}\label{eq:k}
k:=\frac{z^2-3}{2a}\in\Z .
\end{equation}
Indeed $z^2\equiv 3\pmod{2a}$ by \eqref{eq:red} taken modulo $2a$.  Substituting
$z^2=2ak+3$ into \eqref{eq:red} gives the equivalent pair
\begin{equation}\label{eq:pair}
z^2=2ak+3,\qquad b^2=(a-2)^2-(2a-1)k-7 .
\end{equation}
For our point $k=-89278360996331=-13\cdot 431\cdot 7297\cdot 2183641$.

\section{Local conditions}\label{sec:local}

Throughout, $(x,y,z)$ is an integral solution, $x=2a$, $y=b-a$, and $t:=|2a-1|$.

\begin{lemma}[conditions at the primes dividing $a$]\label{lem:a}
$a$ is odd, $z$ is odd, and $z^2\equiv 3\pmod{2a}$.  Every prime $p\ge5$ dividing $a$ satisfies
$p\equiv\pm1\pmod{12}$.  Moreover $9\nmid a$.
\end{lemma}
\begin{proof}
Reduce \eqref{eq:red} modulo $2a$: $-z^2\equiv -3$.  If $a$ were even then $4\mid 2a$ and
$z^2\equiv 3\pmod 4$, which is impossible.  So $z$ is odd.  For $p\ge5$, $p\mid a$, $3$ is a square
mod $p$, i.e.\ $\leg{3}{p}=1$, i.e.\ $p\equiv\pm1\pmod{12}$.  If $9\mid a$ then $3\mid z$ and
$z^2\equiv 3\pmod 9$, which is impossible.
\end{proof}

\begin{lemma}[conditions at $2$]\label{lem:two}
$a\equiv 3\pmod 4$ and $b$ is odd.
\end{lemma}
\begin{proof}
Since $z$ is odd, $z^2\equiv1\pmod 8$, and $8a^2\equiv0$.  Reducing \eqref{eq:red} modulo $8$ gives
$2ab^2+2a-1\equiv 2a^3-3$, i.e.\ $ab^2\equiv a^3-a-1\pmod 4$.  As $a$ is odd, $a^3\equiv a\pmod 8$,
so $ab^2\equiv -1\pmod4$.  This forces $b$ odd and $a\equiv-1\pmod 4$.
\end{proof}

\begin{lemma}[conditions at the primes dividing $2a-1$]\label{lem:t}
$(2b)^2\equiv -19\pmod{t}$.  Consequently every prime $p\mid t$ is either $19$, occurring to the
first power only, or satisfies $\leg{p}{19}=1$.
\end{lemma}
\begin{proof}
Modulo $2a-1$ we have $2a\equiv1$, so \eqref{eq:red} gives
$b^2\equiv N(a)=\tfrac{(2a)^3}{4}-2(2a)^2-3\equiv \tfrac14-2-3=-\tfrac{19}{4}$.  For $p\ne19$,
quadratic reciprocity gives $\leg{-19}{p}=\leg{p}{19}$ because $-19\equiv1\pmod4$.  If $19^2\mid t$,
then $19\mid 2b$ and $19^2\mid(2b)^2+19$, which is impossible.
\end{proof}

\begin{lemma}[the modulus $3$]\label{lem:three}
$a\not\equiv 2\pmod 3$.
\end{lemma}
\begin{proof}
Suppose $a\equiv2\pmod 3$.  Then \eqref{eq:red} modulo $3$ reads $4b^2+3z^2\equiv 16-32-3$, i.e.\
$b^2\equiv 2\pmod3$, which is impossible.
\end{proof}

\begin{proposition}[a reciprocity obstruction]\label{prop:BM}
$a\equiv 3\pmod{36}$, i.e.\ $x\equiv 6\pmod{72}$.
\end{proposition}
\begin{proof}
Consider the quaternion algebra $(x,3)$ evaluated at the solution, i.e.\ the Hilbert symbols
$(x,3)_v$ for all places $v$ of $\Q$.  By the product formula, $\prod_v (x,3)_v=1$.

\emph{Places $v=\infty$ and $p\nmid 6x$.}  Here $(x,3)_v=1$: at $\infty$ because $3>0$, and at $p$
because both entries are units.

\emph{Primes $p\ge5$ with $p\mid x$.}  Here $(x,3)_p=\leg{3}{p}^{v_p(x)}=1$ by
Lemma~\ref{lem:a}.

Hence $(2a,3)_2=(2a,3)_3$.

At $2$: $(2,3)_2=(-1)^{(9-1)/8}=-1$ and $(a,3)_2=(-1)^{\frac{a-1}{2}\cdot\frac{3-1}{2}}=(-1)^{(a-1)/2}=-1$
by Lemma~\ref{lem:two}.  So $(2a,3)_2=+1$.

At $3$:
\begin{itemize}
\item If $3\nmid a$, then $a\equiv1\pmod3$ by Lemma~\ref{lem:three}, and
  $(2a,3)_3=\leg{2a}{3}=\leg{2}{3}=-1\ne 1$.  This is a contradiction, so $3\mid a$.
\item By Lemma~\ref{lem:a}, $a=3a'$ with $3\nmid a'$.  Then
  $(2a,3)_3=(3,3)_3\,(2a',3)_3=\leg{-1}{3}\leg{2a'}{3}=\leg{a'}{3}$.  This must equal $1$, so
  $a'\equiv1\pmod3$, i.e.\ $a\equiv3\pmod 9$.
\end{itemize}
Together with $a\equiv3\pmod4$ this gives $a\equiv3\pmod{36}$.
\end{proof}

\begin{remark}
Proposition~\ref{prop:BM} is an instance of the (integral) Brauer--Manin obstruction.  The class of
$(x,3)$ is unramified on the affine surface away from the boundary.  Locally, $a\equiv1\pmod{12}$
\emph{is} solvable at $2$ and at $3$ separately; it is only the global reciprocity
$\prod_v(x,3)_v=1$ that kills it.  A direct check of $2$-adic and $3$-adic solvability, carried
out by brute force modulo $2^{10}$ and $3^6$, shows that local conditions alone allow
$a\equiv 3\pmod 4$ and $a\equiv 1\pmod 3$ or $a\equiv 3\pmod 9$.  The analogous algebra
$(x-1,-19)$, attached to the other degenerate fibre $2a=1$, turns out to impose nothing new.
\end{remark}

\begin{lemma}[the cubic $N(a)$]\label{lem:N}
Let $p$ be a prime with $p\nmid 2a(2a-1)$ and $p\,\|\,N(a)$.  Then $\leg{-2a(2a-1)}{p}=1$.
\end{lemma}
\begin{proof}
If $p\mid b$ and $p\mid z$ then $p^2\mid N(a)$, contrary to assumption.  If exactly one of $b,z$ is
divisible by $p$, then \eqref{eq:red} modulo $p$ forces $p\mid 2a$ or $p\mid 2a-1$.  So $p\nmid bz$ and
$(z/b)^2\equiv -2a/(2a-1)$, and $-2a(2a-1)$ is a square mod $p$.
\end{proof}

\begin{remark}
One checks that $\gcd(N(a),a)\mid 3$ and $\gcd(N(a),2a-1)\mid 19$.  When $19\mid 2a-1$ there is a
further $19$-adic condition on $a\bmod 19^2$ (in effect $N(a)/19$ must be a square mod $19$), which
the search did not exploit; it would have saved about $1\%$.  For $5\le p\le 23$, $p\ne19$, the conditions above
agree exactly with brute-force $p$-adic solvability modulo $p^3$.
\end{remark}

\subsection*{How restrictive is this?}
Lemmas~\ref{lem:a}--\ref{lem:N} and Proposition~\ref{prop:BM} leave
\begin{itemize}
\item $a$ in one class modulo $36$;
\item $a/3$ composed of primes $\equiv\pm1\pmod{12}$, which is ``half of the primes'';
\item $|2a-1|$ composed of quadratic residues mod $19$ (again half of the primes), plus possibly
  one factor $19$;
\item no small prime of $N(a)$ violating Lemma~\ref{lem:N}.
\end{itemize}
At $|a|\approx 10^{15}$ only about $2\%$ of the integers $a\equiv3\pmod{36}$ (between $1.7\%$ and $2.2\%$ depending on the sign of $a$) survive all four
conditions, i.e.\ about one integer in $1800$.

\subsection*{The found point satisfies every condition}

\begin{center}
\begin{tabular}{>{$}l<{$}>{$}l<{$}}
\toprule
a & -1225206075998589\equiv 3 \pmod{36}\\
a/3 & -\,23\cdot 73\cdot 457\cdot 18803\cdot 28307\quad(\text{residues mod }12:\ 11,1,1,11,11)\\
2a-1 & -\,40161731\cdot 61013609\quad(\text{residues mod }19:\ 6, 11;\ \text{both squares})\\
N(a) & -\,3\cdot 13063\cdot 23567\cdot 1928926393102203289\cdot 2064777611218118087\\
\leg{-2a(2a-1)}{p} & +1 \text{ for each of the four primes } p\ne 3 \text{ of } N(a)\\
a\bmod 8,\ b\bmod 8,\ z\bmod 8 & 3,\ 7,\ 1 \qquad (b,z \text{ odd};\ 3\mid z)\\
\bottomrule
\end{tabular}
\end{center}

\noindent Two of the primes of $N(a)$ have $19$ digits.  No sieve could have certified the condition of
Lemma~\ref{lem:N} for them in advance; the search only ever checks it for $p<1.5\cdot10^6$.

\section{Why the first solution is large}\label{sec:heuristics}

\subsection{Geometry}
Homogenise \eqref{eq:main} with $w$:
$x^2y+xy^2+xz^2+2x^2w-z^2w+3w^3=0$.  This is a smooth cubic surface $S\subset\mathbb P^3$.  Its
section at infinity $\{w=0\}$ is
\[
x\,(xy+y^2+z^2)=0,
\]
the union of the line $L_\infty=\{x=w=0\}$ and the conic $C_\infty=\{xy+y^2+z^2=0=w\}$.  These meet
in the two complex conjugate points $y=\pm iz$.  So the affine surface $U=S\setminus(L_\infty\cup C_\infty)$
is a \emph{log Calabi--Yau surface} whose boundary is a cycle of two rational curves.  The real
points of $U$ go to infinity only along the real circle $C_\infty(\mathbb R)$, which is the circle
$b^2+z^2\approx a^2$ of \S2.2.  For such surfaces integral points are expected to be
sparse, with counts growing only like a power of $\log$ of the height.

The only conic bundle defined over $\Q$ is the pencil of planes through $L_\infty$, i.e.\ the
$x$-fibration of \S2.2, and its fibres are ellipses.  There is therefore no Pell-type family of
solutions; nor did a search find rational lines or twisted cubics on $U$ that could parametrise solutions.

\subsection{A probabilistic count}
Fix an admissible $a$.  The congruence $z^2\equiv3\pmod{2a}$ has $2^{\omega(a/3)}$ roots modulo
$|a|$, where $\omega(n)$ is the number of distinct prime factors of $n$.  Each gives one candidate
$z\in[0,|a|]$ up to sign, and then $b^2$ is determined by \eqref{eq:pair}.  A number of size about
$a^2\cos^2\varphi$ is a perfect square with probability about $1/(2|a\cos\varphi|)$.  The local
conditions modify this: for instance, each prime of $2a-1$ doubles the chance that $b^2$ is a
square.  Summing over admissible $|a|\le A$, the root counts $2^{\omega}$ compensate exactly for
the thinness of the admissible set, and one finds
\[
\#\{\text{solutions with } |a|\le A\}\ \approx\ c\,\log A
\]
with a small constant $c$.  The constant was estimated numerically from the sibling equations
$x^2y+xy^2+xz^2+2x^2-z^2+c=0$ ($c\ne 3$), which have many small solutions.  For $c=3$, the restriction
$a\equiv3\pmod{36}$ and the two quadratic-character conditions make the constant much smaller:
roughly $0.1$--$0.3$ expected solutions per factor $10$ in $|a|$.

The first solution can therefore easily lie at $|a|\approx10^{15}$ or beyond, which is what happened.
The expected number of solutions in any dyadic range $[A,2A]$ is roughly constant, while the cost
of searching it is proportional to $A$.  This is the usual situation for log Calabi--Yau equations,
the sums of three cubes being the best-known example.

\subsection{Algebraic structure: two class-number-one rings}
Two norm forms are hidden in \eqref{eq:pair}.

\emph{The $\Z[\sqrt3]$ side.}  $z^2-3=2a\cdot k$, so $2a$ and $k$ are complementary divisors of the
norm $N_{\Q(\sqrt3)/\Q}(z+\sqrt3)$.  In terms of binary quadratic forms, $[2a,2z,k]$ has
discriminant $12$ and is $\mathrm{SL}_2(\Z)$-equivalent to $X^2-3Y^2$.

\emph{The $\Z[\omega]$ side, $\omega=\frac{1+\sqrt{-19}}2$ (class number $1$).}  Completing the square
in \eqref{eq:pair} gives
\begin{equation}\label{eq:dprime}
(a-k-2-b)(a-k-2+b)=k^2+3k+7=N(k+1+\omega).
\end{equation}
So with $d'=a-k-2-b$ and $e'=a-k-2+b$ we get $d'e'=k^2+3k+7$.  Equivalently, $[e',\,2k+3,\,d']$ is
a binary quadratic form of discriminant $-19$, hence equivalent to $X^2+XY+5Y^2$.  We have
$2a-1=e'+(2k+3)+d'$, and the whole equation collapses to
\begin{equation}\label{eq:topo}
2z^2=(T+1)(T-A-C-3)+6,\qquad A=e',\ C=d',\ T=A+B+C=2a-1,\ B=2k+3 .
\end{equation}
Here $(A,C,T)$ are the values of $X^2+XY+5Y^2$ on the three regions around a vertex of Conway's
topograph, and $B$ is the label on the corresponding edge.  Thus integral points of
\eqref{eq:main} correspond to vertices of the (negative-definite for $a<0$) topograph at which
one quadratic expression is twice a square.

For our point: $d'=-2268341573963755=-5\cdot541\cdot624667\cdot1342433$,
$e'=-3513856040765=-5\cdot7\cdot359\cdot279654281$, $B^2-4AC=-19$, and \eqref{eq:topo} holds.  Note
also $d'=(y^2+y+5)/(x-1)$.

This structure explains Lemma~\ref{lem:t}, since $2a-1$ is a norm from $\Z[\omega]$, and
Lemma~\ref{lem:a}, since $2a$ is a norm from $\Z[\sqrt3]$.  It does not give a shortcut: walking
the topograph costs more per useful candidate than the sieve described below.

\subsection{Caps and elliptic fibrations}\label{sec:caps}
Fixing $d'$ in \eqref{eq:dprime} makes $z^2=k(d'+e')+2k^2+4k+3$, with $e'=(k^2+3k+7)/d'$, a
genus-one curve.  Explicitly it is
\[
E_{d'}:\ W^2=X^3+d'(2d'+3)X^2+d'^2(d'^2+4d'+7)X+3d'^4,\qquad X=d'k,\ W=d'^2z .
\]
Similarly fixing $F=a+2k-2z$ (with $FG=s^2-12$, $s=a-2k$, $G=a+2k+2z$) gives another elliptic
fibration.  Solutions with small $|d'|$ live near $\varphi=0$; solutions with small $|F|$ live
near $\varphi=\pi/2$.  These ``caps'' were searched extensively in exploratory computations:
rank and generator computations for $|d'|,|F|\le 1000$, and divisor sieves for $|k|,|s|\le10^9$
with $|d'|,|F|\le10^6$.
Nothing was found.

The point found here has $d'\approx-2.27\cdot10^{15}$, $F=-2339217298639573$ and
$k\approx-8.9\cdot10^{13}$, all of the same order as $a$.  It is a generic point in the middle
of the surface, far from every cap.  No small-parameter fibre contains it, so the only practical
way to reach it was a complete search over $a$.

\section{The search algorithm}\label{sec:algorithm}

The search enumerates $a=s\,m$ with $s=\pm1$ and $m=r+36j$, where $r=3$ for $s=+1$ and $r=33$ for
$s=-1$, so that $a\equiv3\pmod{36}$ in both cases.  It processes $j$ in segments of $2^{19}$
consecutive values.  The sizes involved are $m$ up to $1.2\cdot10^{15}$ and $|z|\le m+5$.  For each
segment:

\begin{enumerate}
\item \textbf{Exclusion sieve.}  Mark $j$ as impossible when:
  \begin{itemize}
  \item a prime $p\equiv\pm5\pmod{12}$ with $p\le\sqrt{m_{\max}}$ divides $m$ (Lemma~\ref{lem:a});
  \item a prime $p$ with $\leg p{19}=-1$ and $p\le\sqrt{2m_{\max}}$ divides $t=2m-s$
    (Lemma~\ref{lem:t});
  \item $361\mid t$;
  \item $a\equiv\theta\pmod p$ for one of the $57{,}237$ simple roots $\theta$ of $N$ modulo primes
    $p<1.5\cdot10^6$ with $\leg{-2\theta(2\theta-1)}{p}=-1$, unless $a$ lies on the Hensel lift of
    $\theta$ modulo $p^2$ (Lemma~\ref{lem:N}).
  \end{itemize}
  The primes $5,7,13,17,19$ and $29,31,37$ are applied through two precomputed periodic byte
  patterns.  Primes larger than the segment are handled by a bucket sieve, whose entries store
  the prime itself.
\item \textbf{The cofactor of $t$ without factoring.}  After sieving every bad prime up to
  $\sqrt{t}$, at most one prime of $t$ is unknown.  By multiplicativity of the Legendre symbol it
  is good if and only if $t^*$, which is $t$ with one factor $19$ removed, is a square modulo $19$.
  One table lookup per survivor replaces a factorisation of $t$.
\item \textbf{Factor $m/3$} over the good primes $p\equiv\pm1\pmod{12}$, $p\le\sqrt{m_{\max}}$, by a
  second sieve restricted to survivors (a bitmap with rank).  The remaining cofactor $q$ is $1$ or
  a prime, automatically $\equiv\pm1\pmod{12}$.
\item \textbf{Square roots of $3$ modulo $q$}, batched four at a time with Montgomery arithmetic
  and lazy reduction:
  \begin{itemize}
  \item $q\equiv3\pmod4$: $\rho=3^{(q+1)/4}$;
  \item $q\equiv5\pmod8$: Atkin's formula with base $6$;
  \item $q\equiv1\pmod 8$: Tonelli--Shanks, with a non-residue found from quadratic reciprocity
    by table lookups.
  \end{itemize}
  The square roots of $3$ modulo the small primes are precomputed and lifted by Hensel's lemma for
  prime powers.
\item \textbf{CRT with sign symmetry.}  All roots of $z^2\equiv3\pmod m$ are combined from the prime
  roots, with $z\equiv0\pmod 3$.  Only one representative of each pair $\{R,\,m-R\}$ is built.  Both
  odd lifts $z$ and $2m-z$ in $[0,m+5]$ are then tested.
\item \textbf{Candidate tests.}  Compute $k=(z^2-3)/(2a)$ exactly, using an inverse modulo $2^{64}$
  and verifying the result, and
  \[
  B_2=(a-2)^2-(2a-1)k-7=m^2-4sm-z^2+sk .
  \]
  Then apply, in order:
  \begin{itemize}
  \item $B_2>0$;
  \item $B_2\equiv1\pmod8$ (Lemma~\ref{lem:two});
  \item $B_2$ is a square modulo $45045=5\cdot7\cdot9\cdot11\cdot13$ and modulo
    $7429=17\cdot19\cdot23$, using signed residue arithmetic;
  \item an exact integer square root.
  \end{itemize}
  Only about $1.2\%$ of the candidates reach the last step.
\end{enumerate}

Typical statistics for a $2\cdot10^{11}$-wide chunk of $|a|$ near $10^{15}$, for one sign, are
$5.6\cdot10^9$ values of $j$, of which $1.2\cdot10^8$ survive the sieves.  These produce
$3.1\cdot10^8$ candidates $z$, of which $3.7\cdot10^6$ need an exact square root.  The cost is
linear in the range of $|a|$ searched.  Throughput was about $2\cdot10^9$ values of $|a|$ per
second (both signs) on two Intel Xeon (Cascade Lake, 2.8\,GHz) cores, and about
$9\cdot10^9$ per second with eight worker processes on an Apple M1 Pro.

\section{Validation, and a bug worth recording}

\textbf{An independent implementation.}  A short Python/SymPy program, \texttt{check.py}, re-derives
everything from first principles:
\begin{itemize}
\item it factors $a/3$ and $2a-1$;
\item it tests the $N$-condition by trial division over all primes below the bound;
\item it computes all roots of $z^2\equiv3\pmod{2|a|}$ with \texttt{sqrt\_mod};
\item it applies the same candidate tests.
\end{itemize}
On windows at $|a|\approx10^5$, $3.6\cdot10^7$, $3.6\cdot10^{10}$ and $10^{13}$, for both signs, the
C program and \texttt{check.py} agree exactly.  They agree on the counts of kept $a$, of candidates
and of survivors of each filter, and on the full list of $(m,z)$ pairs reaching the exact
square-root test.  Every later optimisation was required to reproduce these lists exactly.  These
were bucket sieving, SIMD scans, lazy Montgomery arithmetic, the portable Mac build and L1 blocking.
The Mac build also ran a self-test of six known windows before starting.

\textbf{The bug.}  An earlier generation of the program computed the residue of $B_2$ modulo
$45045$ and $7429$ with an expression that mixed \texttt{uint64\_t} and \texttt{int64\_t}.  C's usual
arithmetic conversions then evaluated it in unsigned arithmetic.  Whenever the intermediate value
was negative, the residue was taken of $2^{64}+v$ instead of $v$, and a genuine square was rejected
with high probability.  Because all versions shared the line, comparing versions against each
other could not reveal it.  It surfaced when a rewrite produced twice as many
square-root candidates.  The effective coverage of that first run, up to $|a|\approx10^{13}$, was
only about $30\%$.  The whole range was re-searched from zero with the corrected code, which is the
version validated against \texttt{check.py} as above.

\textbf{Sanity checks inside the program.}  The program aborts if any of the following fails:
\begin{itemize}
\item every root satisfies $z^2\equiv3\pmod m$, checked through the exactness of $k$;
\item the CRT modulus equals $m$;
\item the large cofactor is $\equiv\pm1\pmod{12}$.
\end{itemize}
A wrong root therefore cannot silently lose a candidate.

\section{The computation}\label{sec:computation}

\begin{center}
\small
\begin{tabular}{p{0.2\textwidth}p{0.42\textwidth}p{0.3\textwidth}}
\toprule
date and time (CET/CEST) & event & coverage\\\midrule
24 Oct, 00:10 & corrected search restarted from $a=0$ (cloud, 2 cores) & \\
24 Oct, 01:48 & 8 workers started on the M1 Pro & \\
24 Oct, 03:17 & all $|a|<2\cdot10^{13}$ complete & $|x|<4\cdot10^{13}$\\
24 Oct, 03:17 onwards & interleaved chunks of width $2\cdot10^{11}$ & Mac $5/6$, cloud $1/6$\\
25 Oct, 10:03 CET & \textbf{solution} found by Mac worker 2 & chunk $|a|\in[1.2252,1.2254)\cdot10^{15}$\\
\bottomrule
\end{tabular}
\end{center}

The solution came from the $s=-1$ half of the chunk, at $j=34033502111071$, i.e.\ $m=33+36j$ and
$a=-m$.  For this $a$ there were $2^5=32$ roots of $z^2\equiv3\pmod{|a|}$, giving $16$ candidates up
to sign, and exactly one of them made $B_2$ a perfect square.  When the solution appeared, the
workers stopped by themselves.

\textbf{Searched range.}  When the solution appeared, the interleaved layout had left a few chunks
below it unfinished.  These were the cloud's share between $1.18\cdot10^{15}$ and $1.2254\cdot10^{15}$,
plus one interrupted half-chunk: $20$ intervals for $s=+1$ and $21$ for $s=-1$, each $5.6\cdot10^9$
values of $j$ wide.  They were searched afterwards with the same program.
None of the $41$ gap runs found a solution.  The searched range is therefore
complete for both signs of $a$:
\[
|a|<1.226\cdot10^{15},\qquad\text{i.e.}\qquad |x|<2.452\cdot10^{15}.
\]
It was covered by $12{,}425$ completed program runs, with every $j$-interval accounted for in the
logs.  In this whole range the only line ever reported was the one solution above.

\begin{theorem}[computational]\label{thm:unique}
Up to the symmetries $z\mapsto-z$ and $y\mapsto -x-y$, the solution of Theorem~\ref{thm:main} is the
only integral solution of \eqref{eq:main} with $|x|<2.452\cdot10^{15}$.  In particular it is the
solution with the smallest $|x|$.
\end{theorem}

\noindent This statement rests on two things: the proofs in Section~\ref{sec:local}, which justify
every exclusion the program makes, and the correctness of the program.  The program was validated
as described in the previous section.

\section{Remarks}
\begin{itemize}
\item The heuristic suggests infinitely many solutions, with about $c\log A$ of them up to height
  $A$.  The next one is expected to be much further out, by one or more factors of $10$ in $|a|$.
\item The method is a direct analogue of Booker--Sutherland-style searches for sums of three cubes,
  but here every natural factorisation is balanced.  The pairs $2a\cdot k$, $(2a-1)\cdot v$ with
  $v=2a-4k-7$, $d'\cdot e'$ and $F\cdot G$ all have factors of the same size for generic points.
  So there is no ``small factor'' to exploit, and the search must run over $a$ itself.
\item The reciprocity argument of Proposition~\ref{prop:BM} cuts the residue classes of $a$ modulo $36$
  that local conditions allow from four to one.  More refined conditions from the Brauer group, if there are
  any, could reduce it further.
\end{itemize}

\appendix
\section{Reproducing the verification}
\begin{lstlisting}[language=Python]
x, y, z = -2450412151997178, 2357619934960084, 467727250324161
f = lambda x, y, z: x*x*y + x*y*y + x*z*z + 2*x*x - z*z + 3
for Y in (y, -x - y):
    for Z in (z, -z):
        assert f(x, Y, Z) == 0
\end{lstlisting}

\section{Files}
\noindent\begin{tabular}{p{0.42\textwidth}p{0.52\textwidth}}
\texttt{search\_mac.c} & the program that found the solution (portable C, single thread)\\
\texttt{run.sh}, \texttt{worker.sh}, \texttt{workers.sh}, \texttt{config.sh} & 8-worker driver with atomic chunk claims\\
\texttt{selftest.sh}, \texttt{status.sh}, \texttt{stop.sh} & self-test, progress, graceful stop\\
\texttt{search14.c}, \texttt{drive\_cloud.sh} & x86 version (OpenMP, AVX2) and cloud driver\\
\texttt{make\_nbad.py}, \texttt{nbad.txt} & table of excluding roots of $N$ modulo $p<1.5\cdot10^6$\\
\texttt{check.py} & independent Python/SymPy re-implementation used for validation\\
\end{tabular}

\end{document}
